\problemname{Historical Hits}

\newcommand{\hitsmaxn}{3000}
\newcommand{\hitsmaxyear}{10^9}

\emph{Hitster} is a tabletop game about placing songs into an ordered \emph{timeline} by guessing their release year.
As a longtime member of the \emph{Companionless Playing Card Games (CPCG)} association, Harper has access to a single-player variant of Hitster.
The rules are similar to regular Hitster.
Harper starts with an empty timeline.
In each round of the game, she draws a card which contains only a QR code on its front.
She scans the QR code to listen to the song and afterwards guesses the release year of the song.
Afterwards, she turns the card to reveal the song's true release year.
Suppose that the true year is~$a$ and she guessed~$b$.
If Harper's timeline does not yet contain a song with a \emph{true} release year that is strictly between $a$ and $b$, she adds the card to her timeline.\footnote{Her guess $b$ never equals another card's true year, so the card's position in the timeline is uniquely determined.}
Otherwise, she discards the card.

Harper just had a lengthy playthrough of the entire deck of $n$ cards.
Despite her thorough knowledge of music, her final timeline is shorter than she expected.
Playing a second time now would be boring because Harper remembers the correct year of all the songs, along with her not-so-correct guesses.\footnote{She developed this incredible memory through many long nights of playing a single-player variant of \emph{Memory}.}
Surely, she was just unlucky with the order in which she drew the cards, she thinks.
To confirm her suspicion, Harper turns to probability theory.
Suppose that the same cards are presented in an order drawn uniformly at random from the set of all $n!$ permutations.
What is the expected length of Harper's final timeline if she repeats her guesses from her last playthrough for all of the $n$ songs?

\begin{figure}[h]
    \centering
    \includegraphics[width=0.9\textwidth]{sample3}
    \caption{
        Visualization of a playthrough of the game with the cards of Sample 3 presented in the order of the input.
        Harper's timeline is shown at the bottom.
        After the third round, her final timeline contains two cards.
        The top right of each round shows the back of the card, which Harper cannot see before making her guess.
    }
\end{figure}

\begin{Input}
    The input consists of:
    \begin{itemize}
        \item One line with an integer $n$ ($1 \le n \le \hitsmaxn$), the number of cards.
        \item $n$ lines, the $i$th of which contains two integers $a_i, b_i$ ($1 \le a_i, b_i \le \hitsmaxyear$), the correct year and Harper's guessed year for the song of the $i$th card.
    \end{itemize}
    It is guaranteed that $a_i \ne a_j$ and $a_i \ne b_j$ for all $i \ne j$.
\end{Input}

\begin{Output}
    Determine the expected length of Harper's timeline if the cards are presented in an order drawn uniformly at random from the set of all $n!$ permutations.

    Let $M = 998\,244\,353$.
    The expected length can be represented as an irreducible fraction $p / q$ where $q$ is not divisible by $M$.
    Output one integer $p \cdot q^{-1} \bmod M$ (the unique integer $x$ such that $0 \le x < M$ and $x \cdot q \equiv p \bmod M$).
\end{Output}

\section*{Sample Explanation}

In the first sample, there are two possible orders.
If the cards are drawn in the order of the input, then the first card is added to the timeline and the second card is discarded.
If the cards are drawn in reverse order, then both cards are added to the timeline.
Thus, the expected length is $3/2$.
One can verify that $499122178 \cdot 2 \equiv 3 \bmod M$.
